%==============================================================
% BAB IV  HASIL DAN PEMBAHASAN
%==============================================================

\chapter{HASIL DAN PEMBAHASAN}

%--------------------------------------------------------------
\section{Hasil Representasi Softmap Spasial}
%--------------------------------------------------------------

Representasi \textit{softmap} spasial menghasilkan tiga kanal yang saling melengkapi. Kanal
plain Gabor menonjolkan tekstur pembuluh darah yang menyebar dari kutub posterior ke area
perifer retina. Kanal \textit{ridge} berbasis Meijering dan \textit{oriented top-hat}
menonjolkan struktur garis batas serta tonjolan yang relevan untuk membedakan Stage 1 dan
Stage 2. Kanal hijau CLAHE mempertahankan konteks intensitas lokal agar model tetap menerima
informasi jaringan retina yang tidak seluruhnya tercakup oleh dua kanal struktural pertama.

Perbandingan visual pada Gambar \ref{fig:input_scenarios_comparison} menunjukkan bahwa
komposit \textit{softmap} menyajikan informasi yang lebih terstruktur dibanding RGB mentah.
Pada citra Stage 1, kanal \textit{ridge} membantu menonjolkan garis batas yang tipis dan
berkontras rendah. Pada Stage 2 dan Stage 3, kombinasi kanal pembuluh darah dan kanal
\textit{ridge} memberikan petunjuk spasial mengenai penebalan \textit{ridge} dan pertumbuhan
neovaskularisasi. Representasi inilah yang kemudian dibandingkan secara langsung dengan RGB
mentah pada eksperimen klasifikasi.

Pada evaluasi pendukung segmentasi struktural, kanal pembuluh darah mencapai nilai Dice
\textbf{0,3881} pada dataset HVDROPDB-BV. Kanal \textit{ridge} mencapai nilai Dice
\textbf{0,0993} pada HVDROPDB-RIDGE ketika dievaluasi dengan protokol yang sama untuk memeriksa
kesesuaian peta \textit{ridge} dengan anotasi garis batas dan tonjolan retina. Evaluasi ini
berperan sebagai validasi spasial pendukung, sementara eksperimen utama penelitian tetap
berfokus pada dampak representasi \textit{softmap} terhadap performa klasifikasi \textit{stage}
ROP.

\begin{figure}[H]
  \centering
  \includegraphics[width=0.7\linewidth]{images/segmentation_metrics}
  \caption{Ringkasan Dice pendukung untuk kanal pembuluh darah dan kanal \textit{ridge}}
  \label{fig:segmentation_metrics}
\end{figure}

\begin{figure}[H]
  \centering
  \includegraphics[width=\linewidth]{images/segmentation_examples}
  \caption{Contoh visual evaluasi pendukung segmentasi struktural pada anotasi Agrawal2021}
  \label{fig:segmentation_examples}
\end{figure}

%--------------------------------------------------------------
\section{Hasil Baseline Klasik}
%--------------------------------------------------------------

\subsection{Perbandingan Antar Pengklasifikasi}

Baseline klasik menggunakan 48 fitur rekayasa manual dari peta pembuluh darah dan area tepi.
Evaluasi dilakukan menggunakan validasi silang 5-\textit{fold} yang distratifikasi pada 756
citra (kelas Stage 1, Stage 2, Stage 3 dari Zhao2024). Hasil perbandingan antar pengklasifikasi
ditunjukkan pada Tabel \ref{tab:klasik_perbandingan}.

\begin{table}[H]
\centering
\caption{Perbandingan \textit{macro F1-score} antar pengklasifikasi pada baseline klasik
(3 kelas: Stage 1, Stage 2, Stage 3)}
\label{tab:klasik_perbandingan}
\begin{tabular}{l c}
\toprule
\textbf{Pengklasifikasi} & \textbf{\textit{Macro F1}} \\
\midrule
SVM (kernel RBF)           & 0,4964 \\
\textbf{Random Forest}     & \textbf{0,5147} \\
Gradient Boosting          & 0,4912 \\
\midrule
Baseline hanya fitur pembuluh darah & 0,4740 \\
\bottomrule
\end{tabular}
\end{table}

\textit{Random Forest} menghasilkan \textit{macro F1-score} terbaik sebesar 0,5147 dengan
konfigurasi 400 pohon keputusan dan bobot kelas seimbang. Penambahan 12 fitur area tepi
(\textit{ridge}) ke dalam 36 fitur pembuluh darah memberikan peningkatan sebesar +0,041
(dari 0,474 menjadi 0,5147).

\subsection{Pentingnya Fitur Area Tepi}

Analisis tingkat kepentingan fitur (\textit{feature importance}) pada model \textit{Random
Forest} menunjukkan bahwa empat fitur area tepi masuk ke dalam 12 besar fitur terpenting,
seperti ditunjukkan pada Tabel \ref{tab:feature_importance}.

\begin{table}[H]
\centering
\caption{Dua belas fitur terpenting pada model \textit{Random Forest}}
\label{tab:feature_importance}
\begin{tabular}{l c}
\toprule
\textbf{Fitur} & \textbf{Tingkat Kepentingan} \\
\midrule
tortuosity\_proxy         & 0,0427 \\
comp\_size\_std           & 0,0328 \\
vessel\_anisotropy        & 0,0325 \\
int\_r\_std               & 0,0286 \\
\textbf{ridge\_main\_area}     & \textbf{0,0283} \\
\textbf{ridge\_main\_len}      & \textbf{0,0274} \\
\textbf{ridge\_main\_extent}   & \textbf{0,0267} \\
soft\_std                 & 0,0266 \\
comp\_size\_max           & 0,0260 \\
glcm\_correlation\_std    & 0,0258 \\
\textbf{ridge\_resp\_p95}      & \textbf{0,0247} \\
vessel\_density           & 0,0243 \\
\bottomrule
\end{tabular}
\end{table}

\noindent Nilai-nilai ini mengonfirmasi bahwa fitur area tepi (\textit{ridge}) memberikan
kontribusi nyata dalam membedakan \textit{stage} ROP, khususnya antara Stage 2 (yang ditandai
\textit{ridge}) dan Stage lainnya.

\subsection{Batas Atas Performa Klasik}

Meskipun penambahan fitur tepi meningkatkan performa secara keseluruhan, F1-score untuk Stage
1 hampir tidak bergerak dari rentang 0,30--0,31 pada seluruh konfigurasi yang dicoba. Analisis
menunjukkan bahwa filter tepi klasik hanya mampu mendeteksi sekitar 37\% piksel \textit{demarcation
line} yang sebenarnya. Garis batas yang tipis dan berkontras rendah pada Stage 1 merupakan
struktur yang paling sulit dideteksi oleh filter berbasis aturan tanpa kemampuan belajar
adaptif.

%--------------------------------------------------------------
\section{Hasil Klasifikasi dengan Jaringan Saraf Tiruan}
%--------------------------------------------------------------

\subsection{Ringkasan Hasil Eksperimen}

Eksperimen jaringan saraf tiruan dilakukan secara bertahap, dengan setiap konfigurasi
membangun di atas hasil konfigurasi sebelumnya. Tabel \ref{tab:hasil_keseluruhan} merangkum
seluruh hasil yang diperoleh. Seluruh model jaringan saraf dievaluasi menggunakan protokol
\textit{group-aware cross-validation} (\textit{StratifiedGroupKFold} 5-\textit{fold}) yang
diterapkan secara konsisten sejak awal agar seluruh citra dari mata yang sama berada pada
\textit{fold} yang sama (lihat Subbab \ref{subsec:group_aware}).

\begin{table}[H]
\centering
\caption{Ringkasan hasil eksperimen klasifikasi}
\label{tab:hasil_keseluruhan}
\begin{tabular}{l l c l}
\toprule
\textbf{Eksperimen} & \textbf{Input} & \textbf{\textit{Macro F1}} & \textbf{Catatan} \\
\midrule
Baseline klasik terbaik     & 48 fitur    & 0,5147 & RF, 3 kelas \\
CNN RGB (group-aware)       & RGB 3-kanal & 0,6866 & Baseline CNN \\
\textbf{CNN softmap (group-aware)} & \textbf{Softmap} & \textbf{0,7853} & \textbf{Terbaik} \\
CNN + kalibrasi Stage1      & Softmap     & 0,8018 & Kalibrasi pasca-hoc \\
\bottomrule
\end{tabular}
\end{table}

Tabel \ref{tab:hasil_keseluruhan} menunjukkan ringkasan performa dari berbagai konfigurasi eksperimen yang dilakukan. Model baseline klasik menggunakan fitur rekayasa manual menghasilkan performa awal yang rendah (0,5147), yang kemudian berhasil dilampaui oleh model CNN dengan masukan RGB (0,6866). Penggunaan representasi softmap spasial sebagai masukan model TinyResNet memberikan peningkatan performa terbesar hingga mencapai Macro F1-score sebesar 0,7853, menjadikannya sebagai model dasar terbaik dalam penelitian ini.

Untuk mengatasi tingginya prediksi \textit{false positive} pada Stage 1 (di mana presisi awal hanya 0,45 akibat model yang terlalu sensitif memprediksi Stage 1), dilakukan eksperimen tambahan berupa penyesuaian bias logit \textit{post-hoc} (\textit{post-hoc logit-bias adjustment}). Dengan memberikan nilai bias negatif konstan sebesar $-1{,}40$ pada \textit{logit} Stage 1 setelah proses pelatihan selesai, model dapat menyaring prediksi palsu tersebut secara efektif sehingga meningkatkan nilai \textit{Macro F1-score} secara keseluruhan menjadi 0,8018. Meskipun demikian, model dengan softmap dasar tanpa kalibrasi (0,7853) tetap dianggap sebagai model utama yang paling andal untuk generalisasi karena tidak bergantung pada penyesuaian bias pasca-pemrosesan.

\subsection{Protokol Evaluasi \textit{Group-Aware}}
\label{subsec:group_aware}

Dataset ROP yang digunakan memiliki dua karakteristik yang
menuntut pembagian data dilakukan pada tingkat kelompok mata, bukan pada tingkat citra
individual:

\begin{enumerate}
    \item \textbf{Duplikat byte-identik}: Terdapat 12 berkas citra yang persis sama dalam 6
    kelompok, termasuk tiga pasang dengan label kelas yang berbeda (konflik label). Duplikat
    seperti ini dapat muncul di kedua sisi batas pelatihan--pengujian secara bersamaan.
    \item \textbf{Citra satu mata}: Beberapa citra berasal dari mata yang sama tetapi dicatat
    sebagai berkas terpisah. Analisis \textit{normalized cross-correlation} (NCC $\geq$ 0,99)
    mengidentifikasi sekitar 63 pasang citra hampir-duplikat. Pembagian data pada tingkat citra
    akan membiarkan citra dari mata yang sama masuk ke pelatihan dan pengujian sekaligus,
    sehingga model dapat mengidentifikasi mata, bukan penyakitnya.
\end{enumerate}

Untuk mengatasi kedua karakteristik tersebut, seluruh eksperimen klasifikasi jaringan saraf
menggunakan \textit{StratifiedGroupKFold} 5-\textit{fold}, di mana seluruh citra dari satu
kelompok mata ditempatkan pada \textit{fold} yang sama. Protokol ini menghasilkan estimasi
performa yang lebih jujur dan dapat dipercaya, meskipun nilainya lebih rendah
dibanding pembagian tingkat citra.

\subsection{Hasil Terbaik}

Model dengan performa terbaik adalah \textit{TinyResNet} yang dilatih dari awal menggunakan
masukan tiga kanal softmap, dievaluasi dengan \textit{StratifiedGroupKFold} 5-\textit{fold}.
Hasil lengkap ditunjukkan pada Tabel \ref{tab:hasil_terbaik}.

\begin{table}[H]
\centering
\caption{Hasil model terbaik: \textit{TinyResNet} dengan masukan softmap (group-aware CV)}
\label{tab:hasil_terbaik}
\begin{tabular}{l c}
\toprule
\textbf{Metrik} & \textbf{Nilai} \\
\midrule
OOF \textit{Macro F1-score}        & \textbf{0,7853} \\
OOF Akurasi               & 0,8317 \\
OOF Macro Precision       & 0,7867 \\
OOF Macro Recall          & 0,8028 \\
\midrule
\textit{F1} per-\textit{fold} (5 \textit{fold}) & 0,7747 / 0,7881 / 0,7597 / 0,7834 / 0,8203 \\
\bottomrule
\end{tabular}
\end{table}

Performa per-kelas ditunjukkan pada Tabel \ref{tab:performa_per_kelas}.

\begin{table}[H]
\centering
\caption{Performa per kelas model terbaik (OOF, group-aware)}
\label{tab:performa_per_kelas}
\begin{tabular}{l c c c c}
\toprule
\textbf{Kelas} & \textbf{Presisi} & \textbf{Recall} & \textbf{\textit{F1}} & \textbf{Jumlah} \\
\midrule
Normal       & 0,95 & 0,90 & 0,93 & 236 \\
Stage 1      & 0,45 & 0,77 & 0,57 & 94  \\
Stage 2      & 0,70 & 0,61 & 0,65 & 165 \\
Stage 3      & 0,83 & 0,82 & 0,83 & 261 \\
Laser Scars  & 1,00 & 0,92 & 0,96 & 343 \\
\midrule
\textbf{Macro} & \textbf{0,79} & \textbf{0,80} & \textbf{0,79} & \textbf{1.099} \\
\bottomrule
\end{tabular}
\end{table}

Visualisasi matriks konfusi model terbaik ditunjukkan pada Gambar \ref{fig:confusion_matrix}.

\begin{figure}[H]
  \centering
  \includegraphics[width=0.85\linewidth]{images/confusion_matrix_champion}
  \caption{Matriks konfusi TinyResNet pada evaluasi validasi silang group-aware}
  \label{fig:confusion_matrix}
\end{figure}

Untuk menganalisis lebih lanjut keputusan klasifikasi yang dibuat oleh model, beberapa contoh hasil klasifikasi benar (\textit{true prediction}) dan klasifikasi salah (\textit{false prediction}) dari dataset utama disajikan pada Gambar \ref{fig:prediction_examples}.

\begin{figure}[H]
  \centering
  \includegraphics[width=0.85\linewidth]{images/prediction_examples}
  \caption{Contoh visualisasi klasifikasi benar (true prediction) dan salah (false prediction) oleh model TinyResNet}
  \label{fig:prediction_examples}
\end{figure}

Gambar \ref{fig:prediction_examples} menunjukkan bahwa model mampu mengenali dengan tepat kelas Normal (a) karena tidak adanya struktur patologis dan kelas Stage 3 (b) karena adanya neovaskularisasi ekstraretina yang sangat menonjol di kutub posterior. Namun, terdapat kegagalan klasifikasi pada beberapa sampel. Pada Gambar \ref{fig:prediction_examples}(c), citra Stage 2 yang memiliki tonjolan (\textit{ridge}) tipis salah diprediksi sebagai Stage 1. Hal ini disebabkan karena batas visual antara tonjolan tipis (Stage 2 awal) dan garis batas datar (Stage 1 akhir) subjektif secara klinis. Pada Gambar \ref{fig:prediction_examples}(d), citra Stage 1 dengan garis batas kontras rendah salah diklasifikasikan sebagai Normal. Kesalahan ini merupakan yang umum terjadi akibat kemiripan visual yang tinggi.

%--------------------------------------------------------------
\section{Pembahasan}
%--------------------------------------------------------------

\subsection{Representasi Softmap vs. RGB}

Percobaan ablasi yang paling bisa dipercaya dalam penelitian ini adalah perbandingan langsung
antara masukan RGB mentah dengan softmap pada model, data, protokol, dan hyperparameter
yang identik---dengan representasi input adalah satu-satunya variabel yang berbeda. Hasil perbandingan
ditunjukkan pada Gambar \ref{fig:hasil_perbandingan}.

\begin{figure}[H]
  \centering
  \includegraphics[width=0.85\linewidth]{images/results_comparison_chart}
  \caption{Perbandingan macro F1-score antar konfigurasi eksperimen klasifikasi}
  \label{fig:hasil_perbandingan}
\end{figure}

Masukan softmap menghasilkan \textit{macro F1-score} 0,7853 dibanding 0,6866 untuk masukan
RGB pada model yang sama, menunjukkan peningkatan sebesar \textbf{+0,0987}. Peningkatan
ini sepenuhnya dapat dikaitkan dengan perbedaan representasi input, karena tidak ada variabel
lain yang berubah. Ini membuktikan secara langsung bahwa informasi spasial pembuluh darah
dan area tepi yang dikodekan dalam softmap lebih berguna bagi jaringan saraf dibanding nilai
piksel RGB mentah.

\subsection{Hambatan pada Stage 1 dan Stage 2}

Kelas Stage 1 dan Stage 2 tetap menjadi tantangan terbesar dalam klasifikasi. Stage 1 mencapai \textit{F1-score} 0,57 dengan presisi rendah (0,45) namun recall tinggi (0,77), menunjukkan
bahwa model cenderung terlalu sering memprediksi Stage 1. Banyak sampel Stage 2 yang
diprediksi sebagai Stage 1, yaitu kesalahan yang
disebabkan karena Stage 1 dan Stage 2 berbeda hanya dalam tingkat ketebalan \textit{ridge}
yang secara visual sangat mirip (\textit{adjacent-class confusion}).

Stage 2 mencapai \textit{F1-score} 0,65 dengan recall yang lebih rendah (0,61), menunjukkan sebagian
sampel Stage 2 yang bocor ke Stage 1 dan Stage 3. Pola ini konsisten dengan karakteristik penyakit yaitu, Stage 2
sebagai kondisi \textit{transisi} antara Stage 1 dan Stage 3, sehingga batas keputusannya
secara alami lebih tidak jelas.

\subsection{Pendekatan yang Tidak Efektif}

Penelitian ini secara eksplisit mendokumentasikan dua pendekatan yang telah dicoba dan terbukti
tidak efektif:

\textbf{(1) Pendekatan Ordinal Smoothing}: Hipotesis bahwa memperlakukan \textit{stage} sebagai
variabel ordinal (Stage 1 $<$ Stage 2 $<$ Stage 3) dengan menggunakan \textit{ordinal neighbor
smoothing} ($\epsilon = 0{,}15$) akan mengurangi kesalahan antar-\textit{stage} yang
berdekatan. Hasilnya menunjukkan sebaliknya: pada percobaan ablasi yang terkontrol,
penambahan \textit{ordinal smoothing} ke model berbasis softmap dasar justru menurunkan \textit{macro F1} sebesar
0,055. \textit{Recall} Stage 1 yang ditarget untuk ditingkatkan justru turun dari 0,57 menjadi
0,51. Pendekatan ini dinyatakan gagal.

\textbf{(2)} \textit{Oversampling} \textbf{Stage 1}: Hipotesis bahwa memperbanyak sampel Stage 1 (n=94) dalam
setiap \textit{batch} menggunakan \textit{WeightedRandomSampler} dengan faktor $\times$3 akan
memperbaiki \textit{F1-score} Stage 1. Hasilnya, pada dua \textit{fold} yang dijalankan, performa turun dari
0,778 menjadi 0,717 dan dari 0,773 menjadi 0,684. Masalah pada Stage 1 bukan kekurangan
eksposur (\textit{recall} sudah 0,77), melainkan terlalu banyak prediksi \textit{false positive}
(\textit{precision} hanya 0,45). \textit{Oversampling} memperburuk ketidakseimbangan ini. Pendekatan
ini dinyatakan gagal.

Pelajaran dari kedua kegagalan ini adalah bahwa masalah Stage 1--Stage 2 merupakan konfusi
yang disebabkan oleh kesamaan visual antar kelas yang berdekatan secara klinis. Solusi yang
lebih baik adalah pendekatan yang dapat mengurangi prediksi \textit{false positive} Stage 1 secara
spesifik, bukan pendekatan yang memperbesar sinyal Stage 1 secara global.

\subsection{Dampak Praktis dan Relevansi Klinis}

Hasil penelitian ini memiliki relevansi klinis yang signifikan untuk mendukung program skrining ROP di unit perawatan intensif bayi prematur (NICU). Dalam praktik medis, beberapa dampak praktis dari sistem yang diusulkan adalah sebagai berikut:

\begin{enumerate}
    \item \textbf{Alur Kerja Triage yang Aman dan Efektif}: Dari perspektif penapisan (skrining), sensitivitas tinggi pada kelas berisiko adalah prioritas utama. Model terbaik kami mencapai \textit{recall} sebesar 0,77 pada Stage 1 dan 0,82 pada Stage 3. Hal ini berarti model berperilaku sangat konservatif dan aman secara medis, meminimalkan risiko bayi dengan ROP aktif lolos dari pemeriksaan (false negatives).
    \item \textbf{Membantu Pengambilan Keputusan Rujukan}: Kemampuan klasifikasi yang kuat pada kelas Normal (\textit{F1-score} 0,93) dan kelas \textit{laser scars} (\textit{F1-score} 0,96) memungkinkan model berfungsi sebagai alat penyaring awal otomatis. Bayi dengan kondisi normal dapat disaring dengan cepat, sehingga dokter spesialis mata dapat mencurahkan waktu dan sumber daya medis yang terbatas untuk berfokus pada bayi dengan gejala Stage 2 dan Stage 3 yang membutuhkan pemantauan ketat serta tindakan klinis segera.
    \item \textbf{Mengurangi Variabilitas Antar-Dokter}: Penilaian tingkatan ROP secara visual oleh manusia bersifat subjektif, terutama pada batas transisi antara Stage 1 dan Stage 2. Model komputer memberikan metrik analisis spasial yang objektif dan konsisten yang mengurangi inkonsistensi diagnosis antar-pemeriksa di rumah sakit yang berbeda.
\end{enumerate}
Secara keseluruhan, pemanfaatan representasi komposit spasial dalam arsitektur TinyResNet terbukti memberikan hasil diagnosis yang andal dari segi akurasi dan klinis tanpa memerlukan \textit{transfer learning} dari domain non-medis.

\subsection{Keterbatasan Penelitian}

Beberapa keterbatasan penting perlu dicatat:

\begin{itemize}
    \item \textbf{Tidak ada identitas pasien yang terverifikasi}: Pengelompokan berdasarkan NCC
    merupakan pendekatan terbaik yang tersedia mengingat tidak adanya ID pasien dalam nama
    berkas. Meskipun ini mengurangi risiko tumpang tindih antar himpunan data, hasilnya tidak
    dapat dijamin setara dengan pembagian berbasis pasien yang terverifikasi.
    \item \textbf{Dataset tunggal}: Seluruh eksperimen dilakukan pada dataset Zhao2024 dari
    satu negara dan sejenis kamera. Kemampuan generalisasi ke dataset dari kamera
    berbeda atau populasi yang berbeda belum diuji.
    \item \textbf{Tidak ada bobot pra-latih}: Keputusan untuk tidak menggunakan bobot pra-latih
    (karena kanal masukan adalah softmap, bukan RGB) berarti penelitian ini tidak dapat
    memanfaatkan pengetahuan visual umum yang telah dipelajari dari dataset berskala besar.
    Ini adalah \textit{trade-off} yang disengaja dan terdokumentasi.
    \item \textbf{Batas atas Stage 1}: F1-score Stage 1 sebesar 0,57 masih jauh di bawah
    kelas lainnya. Meningkatkan performa Stage 1 memerlukan penelitian lebih lanjut, mungkin
    dengan menggunakan model segmentasi yang dapat dipelajari (\textit{learned segmenter})
    seperti U-Net \citep{Ronneberger2015} untuk menghasilkan peta tepi yang lebih akurat.
\end{itemize}
